\chapter{The Register Interface and Regtool}
\label{ch:regtool}

\chapabstract{The CSR half of an accelerator: the \code{register\_interface}
bus that carries control/status/configuration, and \tool{regtool}
(``reggen''), which turns one \file{.hjson} description into the RTL and the
C header for it. \code{make regs} is the Makefile target; this chapter is
what it actually does.}

\section{Two buses, one job each}

\file{cordic\_accel.sv} (\cref{ch:fusesoc}) has two interfaces, and they are
not interchangeable: an \textbf{OBI} port moves the angles and results (the
data), and a \textbf{register interface} port moves configuration -- write
\code{SRC}, \code{DST}, \code{N}, write \code{CTRL.START}, poll
\code{STATUS.DONE}. Splitting them is deliberate: the CPU only ever touches
a handful of small registers, and the accelerator moves the bulk data
itself over OBI, which is where the actual speedup comes from.

\code{register\_interface} (pulp-platform, vendored the same way as
\cref{ch:vendoring} describes) defines that second bus: a request struct
(address, write, write data, byte strobes) and a response struct (read
data, error), typedef'd for one specific address/data width with the same
kind of macro \cref{ch:fusesoc} showed for OBI:

\begin{svcode}[caption={rtl/cordic\_accel\_types\_pkg.sv -- the reg bus types}]
typedef logic [31:0] reg_addr_t;
typedef logic [31:0] reg_data_t;
typedef logic [ 3:0] reg_strb_t;

`REG_BUS_TYPEDEF_ALL(cordic_reg, reg_addr_t, reg_data_t, reg_strb_t)
\end{svcode}

\begin{ruleofthumb}
You will not hand-write anything that speaks this protocol directly. You
describe your registers in \file{.hjson}, \tool{regtool} generates a block
that speaks \code{reg\_req}/\code{reg\_rsp} on one side, and your own
design only ever touches plain struct fields on the other.
\end{ruleofthumb}

\section{Anatomy of the .hjson register description}

\begin{corecode}[caption={data/cordic\_accel.hjson (trimmed to two registers)},label={lst:regtool:hjson}]
{ name: "cordic_accel"
  clock_primary: "clk_i"
  reset_primary: "rst_ni"
  bus_interfaces: [
    { protocol: "reg_iface", direction: "device" }
  ]
  regwidth: 32
  registers: [

    { name:     "CTRL"
      desc:     "Control register. Write START=1 to launch a run."
      swaccess: "wo"
      hwaccess: "hro"
      hwqe:     "true"
      fields: [
        { bits: "0", name: "START"
          desc: "Write 1 to start processing N angles."
        }
      ]
    }

    { name:     "STATUS"
      desc:     "Status register, driven by hardware."
      swaccess: "ro"
      hwaccess: "hwo"
      hwext:    "true"
      fields: [
        { bits: "0", name: "DONE" desc: "Set when results are ready." }
        { bits: "1", name: "BUSY" desc: "Set while the port is owned." }
      ]
    }
  ]
}
\end{corecode}

The header names the block and its clock/reset, and \kw{regwidth} fixes
every register at the same width (32, here). Offsets are not written down
anywhere: \tool{regtool} assigns them in declaration order, word aligned --
\code{CTRL} at \code{0x00}, \code{STATUS} at \code{0x04}, and so on. Each
register then has:

\begin{itemize}
  \item \kw{swaccess}: what \emph{software} may do -- \code{wo}
        (write-only), \code{ro} (read-only) or \code{rw}.
  \item \kw{hwaccess}: what \emph{hardware} may do -- \code{hro} (hardware
        reads the value, never writes it: the normal case for a
        configuration register like \code{CTRL} or \code{N}) or \code{hwo}
        (hardware writes it, software only reads: the normal case for a
        status register). Consult \tool{regtool}'s own documentation for
        the remaining combinations; these two cover everything in this
        example.
  \item \kw{resval}: the reset value, when the register needs one other
        than zero (\cref{lst:regtool:hjson} omits it here, but
        \code{SEED\_X} in the full file resets to the CORDIC gain constant).
  \item \kw{fields}: one or more named bit ranges inside the register,
        each with its own \kw{bits} and \kw{desc}.
\end{itemize}

\begin{gotcha}{\code{hwqe} and \code{hwext} look similar and do opposite jobs}
\kw{hwqe: true} adds a one-cycle write-strobe your hardware can see
(\code{reg2hw.ctrl.qe}), on top of the stored value (\code{reg2hw.ctrl.q}).
Use it when hardware needs to know \emph{that a write just happened}, not
only the value -- \cref{sec:regtool:wiring} shows exactly this, detecting a
write of \code{START=1}.

\kw{hwext: true} removes the register's own storage: a software read
returns whatever \code{hw2reg} is driving \emph{this cycle}, combinationally,
instead of a value latched on some earlier cycle. \code{STATUS} needs this
because, without it, software that writes \code{START} and immediately
reads \code{STATUS} would see the flopped value from before the write --
the previous run's \code{DONE} -- and wrongly conclude the new run had
already finished.
\end{gotcha}

\section{What make regs actually runs}

\begin{shcode}
python3 vendor/pulp_platform/register_interface/vendor/lowrisc_opentitan/util/regtool.py \
  -r -t rtl data/cordic_accel.hjson
python3 .../regtool.py -D -o sw/cordic_accel_regs.h data/cordic_accel.hjson
\end{shcode}

\tool{regtool} (``reggen'') is OpenTitan's register-generation tool. It
ships vendored inside \code{register\_interface} (\cref{ch:vendoring}
fetched it along with everything else in that IP), so it is present the
moment \code{make vendor} has run; nothing extra to install. Two runs, two
outputs:

\begin{itemize}
  \item \code{-r -t rtl}: writes \file{rtl/cordic\_accel\_reg\_pkg.sv} (the
        struct \emph{types}) and \file{rtl/cordic\_accel\_reg\_top.sv} (the
        module that speaks the bus and exposes those structs).
  \item \code{-D -o <file>}: writes a C header of offsets, bit positions and
        masks -- one \code{\#define} block per register, generated straight
        from the same \kw{desc} strings, so the comments in
        \file{cordic\_accel\_regs.h} and in the \file{.hjson} match by
        construction.
\end{itemize}

The lab Makefile wraps both calls as \code{make regs}
(\cref{ch:makefiles}); re-run it whenever \file{data/cordic\_accel.hjson}
changes, the same way you re-run \code{make vectors} after changing the
golden model.

\begin{gotcha}{The generated files are not yours to edit}
\file{cordic\_accel\_reg\_pkg.sv}, \file{cordic\_accel\_reg\_top.sv} and
\file{cordic\_accel\_regs.h} all say ``auto-generated'' at the top. Editing
them by hand works until the next \code{make regs} silently overwrites your
change. Edit \file{data/cordic\_accel.hjson} and regenerate instead.
\end{gotcha}

\section{Reading the generated types}

\tool{regtool} turns every register into a small packed struct, and groups
them into one \code{reg2hw\_t} (hardware reads this) and one
\code{hw2reg\_t} (hardware drives this):

\begin{svcode}[caption={rtl/cordic\_accel\_reg\_pkg.sv (generated, excerpt)}]
typedef struct packed {
  logic q;   // the stored value
  logic qe;  // one-cycle pulse: software wrote this cycle
} cordic_accel_reg2hw_ctrl_reg_t;

typedef struct packed {
  struct packed { logic d; } done;  // hardware drives this bit
  struct packed { logic d; } busy;
} cordic_accel_hw2reg_status_reg_t;
\end{svcode}

The naming is consistent everywhere this pattern is used, in every X-HEEP
peripheral and not only this one: a \kw{reg2hw} field ends in \code{.q}
(the value) and, if \kw{hwqe} was set, \code{.qe} (the write pulse); a
\kw{hw2reg} field ends in \code{.d} (the value your hardware should drive).
Once you can read this struct, you can read the CSR block of any peripheral
built the same way.

\section{Wiring it into your design}
\label{sec:regtool:wiring}

\file{cordic\_accel.sv} instantiates the generated \code{reg\_top} once and
never touches the bus signals again:

\begin{svcode}[caption={rtl/cordic\_accel.sv -- instantiating the CSR block}]
cordic_accel_reg2hw_t reg2hw;
cordic_accel_hw2reg_t hw2reg;

cordic_accel_reg_top #(
  .reg_req_t (reg_req_t),
  .reg_rsp_t (reg_rsp_t)
) u_reg_top (
  .clk_i, .rst_ni,
  .reg_req_i,
  .reg_rsp_o,
  .reg2hw,
  .hw2reg,
  .devmode_i (1'b1)
);

// Software wrote START=1 this cycle -> launch a run.
assign start = reg2hw.ctrl.qe & reg2hw.ctrl.q;
...
// Drive STATUS back to software.
assign hw2reg.status.done.d = done_q & ~run_start;
assign hw2reg.status.busy.d = busy_q  | run_start;
\end{svcode}

From here on the rest of the module is ordinary SystemVerilog: an FSM reads
\code{reg2hw.n.q}, \code{reg2hw.src.q}, \code{reg2hw.dst.q} to know what to
fetch and where to write it, and drives \code{hw2reg.status} combinationally
so a read of \code{STATUS} is never stale. Nothing downstream of
\code{u\_reg\_top} knows or cares that a register bus exists at all.

\section{The C side}

The generated header gives software plain integer offsets and bitfield
helpers, nothing more:

\begin{cppcode}[caption={sw/cordic\_accel\_regs.h (generated, excerpt)}]
#define CORDIC_ACCEL_CTRL_REG_OFFSET 0x0
#define CORDIC_ACCEL_CTRL_START_BIT 0

#define CORDIC_ACCEL_STATUS_REG_OFFSET 0x4
#define CORDIC_ACCEL_STATUS_DONE_BIT 0
#define CORDIC_ACCEL_STATUS_BUSY_BIT 1
\end{cppcode}

Writing an actual driver around these offsets is a Lab 3 problem (the
accelerator has to be mapped into X-HEEP's address space first); this
booklet stops at ``where these numbers come from''.

\section{Checklist}
\label{sec:regtool:checklist}

\begin{itemize}
  \item The register interface carries configuration; OBI (or whichever
        data bus a design uses) carries data. Do not mix the two.
  \item One \file{.hjson} per block; \kw{swaccess}/\kw{hwaccess} say who
        may read or write each register, \kw{hwqe} adds a write-pulse,
        \kw{hwext} removes the register's own storage.
  \item \code{make regs} runs \tool{regtool} twice: once for the RTL
        (\code{reg\_pkg.sv} + \code{reg\_top.sv}), once for the C header.
  \item \kw{reg2hw.<reg>.q}/\code{.qe} is what hardware reads;
        \kw{hw2reg.<reg>.d} is what hardware drives. Every peripheral built
        this way uses the same two suffixes.
  \item Never hand-edit a generated file; edit the \file{.hjson} and
        re-run \code{make regs}.
\end{itemize}
